\documentclass{article}
\usepackage[T1]{fontenc}
\usepackage{lmodern}
\usepackage{graphicx}
\usepackage{amsmath,amssymb}
\usepackage[a4paper,margin=2cm]{geometry}
\usepackage[absolute,overlay]{textpos}
\setlength{\TPHorizModule}{1mm}
\setlength{\TPVertModule}{1mm}
\usepackage[indonesian]{babel}
\usepackage{hyperref}
\setlength{\emergencystretch}{2em}
% unit: Halaman 16

\begin{document}

% === Halaman 16 (llm) ===

\begin{flushright}\fbox{Dekripsi: $m = c^d \pmod{n}$}\end{flushright}

\begin{itemize}
    \item Alice mendekripsi cipherteks $C = 0328, 0301, 2653, 2986, 1164$ dengan menggunakan kunci privatnya, yaitu $d = 1019, n = 3337$:
    
    \begin{align*}
    c_1 = 0328 \rightarrow m_1 &= 328^{1019} \pmod{3337} = 704 = 0704 \\
    c_2 = 0301 \rightarrow m_2 &= 301^{1019} \pmod{3337} = 1111 \\
    c_3 = 2653 \rightarrow m_3 &= 2653^{1019} \pmod{3337} = 1400 \\
    c_4 = 2986 \rightarrow m_4 &= 2986^{1019} \pmod{3337} = 1108 \\
    c_5 = 1164 \rightarrow m_5 &= 1164^{1019} \pmod{3337} = 204
    \end{align*}
    
    \item Alice memperoleh kembali plainteks dari Bob \\ $m = 07041111140011080204$
\end{itemize}

yang dikodekan kembali (A = 00, B = 01, C = 02, ..., Z = 25) menjadi \\ $M = \text{HELLO ALICE}$

\end{document}
